\centerline{\bf \Huge Лемма о венике}

\begin{flushright}
        \scriptsize{}
\end{flushright}

\problem{1}
\textbf{Лемма о трезубце.} Пусть $I$ и $I_A$ - центры вписанной и вневписанной (касающейся стороны $B C$ и продолжений сторон $A B$ и $A C$ ) окружностей треугольника $A B C$ соответственно. Тогда точки $B, C, I, I_A$ лежат на одной окружности с центром в середине «меньшей» дуги $B C$ описанной окружности треугольника $A B C$.


\problem{2}
\textbf{Внешняя лемма о трезубце.} Пусть $I_B$ и $I_C$ - центры вневписанных окружностей треугольника $A B C$, касающихся сторон $A C$ и $A B$ соответственно. Тогда точки $B, C, I_C, I_B$ лежат на одной окружности с центром в середине «большей» дуги $B C$ описанной окружности треугольника $A B C$.


\problem{3}
Отрезок, соединяющий середины «меньших» дуг $A B$ и $A C$ описанной окружности треугольника $A B C$, пересекает стороны $A B$ и $A C$ в точках $P$ и $Q$. Докажите, что $A P I Q$ - ромб, где $I$ - центр вписанной в треугольник $A B C$ окружности.


\problem{4}
Вокруг прямоугольного треугольника $ABC$ с прямым углом $C$ описана окружность, на меньших дугах $AC$ и $BC$   взяты их середины --- $K$ и $P$  соответственно. Отрезок $KP$ пересекает катет $AC$ в точке $N$.   Центр вписанной окружности треугольника $ABC$  ---  $I$.   Найти угол $NIC$.


\problem{5}
Биссектрисы углов $B$ и $C$ треугольника $A B C$ пересекают его описанную окружность в точках $B_0$ и $C_0$ соответственно и пересекаются в точке $I$. Окружности $\omega_B$ и $\omega_C$ с центрами $B_0$ и $C_0$ касаются отрезков $A C$ и $A B$ соответственно. Докажите, что через точку $A$ можно провести прямую, касающуюся обеих окружностей $\omega_B$ и $\omega_C$.


\problem{6}
\textbf{Окружность Эйлера.} Пусть $A B C$ - некоторый треугольник. Тогда на одной окружности лежат следующие девять точек: середины $A_1, B_1, C_1$ сторон $B C, C A, A B$ соответственно, основания высот $A_2, B_2, C_2$, опущенных из $A$, $B, C$, середины $A_3, B_3, C_3$ отрезков $A H, B H, C H$, где $H$ - ортоцентр треугольника $A B C$.

\problem{7}
\textbf{Japanese theorem for cyclic quadrilaterals.}
В окружность вписан четырёхугольник $ABCD$ Отметили центры окружностей, вписанных в треугольники $BCD$, $CDA$, $DAB$, $ABC$. Докажите, что отмеченные точки являются вершинами прямоугольника.